\documentclass[a4paper,10pt,twoside,openany]{article}

\usepackage[lang=hebrew]{maths}
\usepackage{hebrewdoc}
\usepackage{stylish}
\usepackage{lipsum}
\let\bs\blacksquare

\usepackage{siunitx,array}
\sisetup{table-format=-1.5, group-digits=false}
\newcolumntype{L}{>{\phantom{$-$}}l}       % prefix some whitespace
\newcommand\mcL[1]{\multicolumn{1}{L}{#1}} % handy shortcut macro

\setlength{\parindent}{0pt}

%%%%%%%%%%%%
% Styling %
%%%%%%%%%%%%

\usepackage{enumitem}

%%%%%%%%%%%%%
% Counters  %
%%%%%%%%%%%%%

\setcounter{section}{1}     
            
%BIBLIOGRAPHY
\usepackage[
backend=biber,
style=alphabetic,
]{biblatex}
\addbibresource{bibliography.bib} %Imports bibliography file

%%%%%%%%%%
% Title  %
%%%%%%%%%%
\title{
אלגברה ב' - גיליון תרגילי בית 6 \\
הפולינום המינימלי, מרחבים שמורים, ואופרטורים נילפוטנטיים
\\
\vspace{1cm}
\large{תאריך הגשה: 14.6.2026}
}
\date{}

\begin{document}
\maketitle

\begin{exercise}
יהי
$V = \CC_3\brs{x}$
ויהי
\begin{align*}
D \colon V &\to V \\
\text{.} \hphantom{l} p\prs{x} &\mapsto p\prs{x} + p'\prs{x}
\end{align*}

מיצאו את הפולינומים האופייני והמינימלי של כל אחד מהאופרטורים הבאים.

\begin{enumerate}
\item $D$
\item $D + \id_V$
\item $D^2$
\item $D^{-1}$
\end{enumerate}
\end{exercise}

\begin{exercise}
יהי
$V$
מרחב וקטורי סוף־מימדי מעל שדה
$\FF$,
ויהי
$T \in \End_\FF\prs{V}$.
\begin{enumerate}
\item תהי
$P \in \End_\FF\prs{V} \setminus \set{0_V, \id_V}$
הטלה עבורה
$TP = PT$.
הראו כי
$\im\prs{P}, \ker\prs{P}$
הינם
$T$%
־שמורים והסיקו כי
$T$
אופרטור פריד.

\item נניח כי
$T$
אופרטור פריד. מיצאו הטלה
$P \in \End_\FF\prs{V} \setminus \set{0_V, \id_V}$
עבורה
$TP = PT$.

\item הסיקו כי $T$ הינו פריד אם ורק אם קיימת הטלה
$P \in \End_\FF\prs{V} \setminus \set{0_V, \id_V}$
המתחלפת עם
$T$.
\end{enumerate}
\end{exercise}

\begin{exercise}
יהי
$V$
מרחב וקטורי סוף־מימדי מעל
$\CC$.
נסמן ב־%
$\mathrm{Nil}\prs{V}$
את אוסף האופרטורים הנלפוטנטיים על
$V$.

\begin{enumerate}
\item מיצאו אופרטורים
$S,T \in \mathrm{Nil}\prs{V}$
עבורם
$ST \notin \mathrm{Nil}\prs{V}$.

\item הראו כי אם
$S,T \in \mathrm{Nil}\prs{V}$
מתחלפים אז
$ST \in \mathrm{Nil}\prs{V}$.

\item מיצאו אופרטורים
$S,T \in \mathrm{Nil}\prs{V}$
עבורם
$S+T \notin \mathrm{Nil}\prs{V}$.

\item הראו כי אם
$S,T \in \mathrm{Nil}\prs{V}$
מתחלפים אז
$S+T \in \mathrm{Nil}\prs{V}$.
\end{enumerate}
\end{exercise}

\begin{exercise}
יהי
$V$
מרחב וקטורי מעל שדה
$\FF$
ויהי
$T \in \End_\FF\prs{V}$
נילפוטנטי.

\begin{enumerate}
\item הוכיחו כי
$\id_V - T$
הפיך ומיצאו את ההופכי שלו.

\textbf{רמז:}
מיצאו פולינום
$p\prs{x}$
הקשור לאופרטור $\id_V - T$, חישבו על דרך אחרת לכתוב את הביטוי
$\frac{1}{p\prs{x}}$
והציבו
$T$
בביטוי שקיבלתן.

\item הוכיחו כי
$\id_V + T$
הפיך והיעזרו בסעיף הקודם כדי למצוא את
$\prs{\id_V + T}^{-1}$.
\end{enumerate}
\end{exercise}

\end{document}